\begin{frame}{Q3 — Prefix caching과 prefill 버킷 설계}
  \hd{The effect}
  \begin{itemize}
    \item Without caching, buckets are sized to the full prompt length.
    \item With a hit, tokens actually computed $= R = (\text{prompt}) - (\text{cached prefix})$;
          $R$ varies per request, usually small.
    \item Pick the bucket from $R$, not the prompt — a prompt-sized graph pads away the benefit.
  \end{itemize}
  \hd{Approach}
  \begin{itemize}
    \item Look up cache $\to$ compute $R$ $\to$ smallest compiled bucket $\ge R$.
    \item One chunk-sized prefill graph $+$ a few tail buckets; long $R$ repeats the chunk.
    \item Bucket sizes from the real post-cache $R$ distribution, recompiled offline.
    \item vllm-rbln: one prefill graph at chunk size; sub-block reuse; matched KV copied before the forward.
  \end{itemize}
\end{frame}

\begin{frame}{Q3 — how prefix caching moves the workload}
  \vspace{1.2em}
  \begin{center}\resizebox{\textwidth}{!}{% S = cached prefix + residual R  ->  bucket is chosen from R
\begin{tikzpicture}[font=\footnotesize, x=1cm, y=1cm]
  % Row 1: no caching
  \node[anchor=east, brick, font=\footnotesize\bfseries] at (-0.2,1.35) {No caching};
  \filldraw[fill=soft, draw=softtxt] (0,1.05) rectangle (7.4,1.65);
  \node[softtxt] at (3.7,1.35) {prompt length $S$ — compute all of it};
  \draw[-{Stealth}, softtxt, thick] (7.7,1.35) -- (8.4,1.35);
  \node[anchor=west, softtxt, font=\footnotesize\bfseries] at (8.55,1.35) {bucket $=S$};

  % Row 2: cache hit
  \node[anchor=east, brick, font=\footnotesize\bfseries] at (-0.2,0.0) {Cache hit};
  \filldraw[fill=soft, draw=softtxt] (0,-0.3) rectangle (5.2,0.3);
  \node[softtxt] at (2.6,0.0) {cached prefix · KV reused};
  \filldraw[fill=brick, draw=brick] (5.2,-0.3) rectangle (7.4,0.3);
  \node[white, font=\footnotesize\bfseries] at (6.3,0.0) {$R$};
  \node[softtxt, font=\scriptsize] at (6.3,-0.62) {varies per request};
  \draw[-{Stealth}, brick, thick] (7.7,0.0) -- (8.4,0.0);
  \node[anchor=west, brick, font=\footnotesize\bfseries] at (8.55,0.0) {bucket $=R$};
\end{tikzpicture}
}\end{center}
  \vspace{1.6em}
  \footnotesize A short $R$ still attends over the \emph{full} cached context
  $\Rightarrow$ the graph needs a small query bucket \textbf{and} a bounded pointer to the
  cached KV (block table $+$ seq-lens), not one large shape.
\end{frame}
